% ==================== 2.3.2 小结表 ====================
表~\ref{tab:summary} 汇总了本文建模的八种常见晶体结构的空间群、晶格常数、配位数与
原子堆积因子。

\begin{table}[H]
  \centering
  \small
  \caption{八种常见晶体结构的结构特征汇总}
  \label{tab:summary}
  \begin{tabular}{@{}llllll@{}}
    \toprule
    结构 & 空间群 & 晶格常数（\AA） & $Z$ & 配位数 & 堆积因子 \\
    \midrule
    FCC（Fe） & $F\,m\,\bar3\,m$ & $a=3.43$ & 4 & 12 & 0.7405 \\
    BCC（Fe） & $I\,m\,\bar3\,m$ & $a=2.819$ & 2 & 8（次近邻 6） & 0.6802 \\
    HCP（Mg） & $P\,6_3/m\,m\,c$ & $a=3.2085$，$c=5.2106$ & 2 & 12 & 0.7405（理想） \\
    NaCl & $F\,m\,\bar3\,m$ & $a=5.62$ & 4 & 6 & — \\
    金红石（TiO$_2$） & $P\,4_2/m\,n\,m$ & $a=4.6$，$c=2.9634$ & 2 & 6 & — \\
    锐钛矿（TiO$_2$） & $I\,4_1/a\,m\,d$ & $a=3.784$，$c=9.5146$ & 4 & 6 & — \\
    萤石（CaF$_2$） & $F\,m\,\bar3\,m$ & $a=5.462$ & 4 & Ca：8，F：4 & — \\
    金刚石（Si） & $F\,d\,\bar3\,m$ & $a=5.42$ & 8 & 4 & 0.3401 \\
    \bottomrule
  \end{tabular}
\end{table}

可见：金属单质以 FCC、BCC、HCP 三种密堆积（或近密堆积）结构为主；离子化合物中，
阴阳离子半径比决定了配位多面体类型（NaCl 六配位、萤石中 Ca 八配位）；共价晶体
（金刚石）以 $sp^3$ 四面体配位为特征，堆积因子最低。晶体结构特征可概括为：
``空间群 + 晶格常数 + 原子占位 $\Rightarrow$ 配位多面体 $\Rightarrow$ 堆积方式''。
